\section{Problem setting and propagation graphs}\label{sec:setting}

A landscape consists of a finite set $V$ of cells and a set $C\subseteq V$ of
cells eligible for treatment. Cell $k$ has a nonnegative loss weight $w_k$;
when all weights equal one, total loss counts burned cells. Treatment is a
first-stage decision: $y_i=1$ means that eligible cell $i$ receives a
firebreak before the fire realization is known. The treatment budget is
$B=\lfloor\alpha |V|\rfloor$, where $\alpha$ is the allowed fraction of
landscape cells. Throughout, placements belong to
\begin{equation}\label{eq:placements}
 \mathcal Y=\left\{y\in\{0,1\}^{C}:\sum_{i\in C}y_i\le B\right\}.
\end{equation}
An eligible ignition cell may be selected, but treatment does not prevent
ignition there: the ignition burns and can transmit fire. Selecting a
non-ignition cell blocks propagation \emph{into} that cell. This convention
is important both for the optimization model and for evaluating a fixed
placement.

We represent ignition and propagation uncertainty by a finite sample
$\mathcal S$ of scenarios, with probabilities $\rho_s>0$ summing to one.
Cell2Fire provides spatial fire-spread realizations
\citep{pais2021cell2fire}. For scenario $s$, let $r_s$ be its ignition
and $G_s^R=(V_s,E_s^R)$ the original directed \emph{Reburn} graph. An arc
$(i,j)$ records a possible propagation from $i$ to $j$ in this recorded
realization. Multiple directed routes can reach the same cell. We take
$V_s$ to consist of the cells reachable from $r_s$ before treatment; other
landscape cells contribute zero to this scenario's loss.

Our optimization representation is a rooted shortest-path spanning tree
$G_s^T=(V_s,E_s^T)$ extracted from $G_s^R$. The tree has the same ignition
and cells, but selects one directed ignition-to-cell route, so that
$E_s^T\subseteq E_s^R$ and every $k\in V_s$ has a unique path from $r_s$.
Arc costs and tie breaking used to construct a particular data set do not
enter the results below; its essential properties are spanning, rooting,
and arc inclusion. Retaining only one route reduces the propagation
structure and permits efficient graph-specific bounds. It also removes
alternative routes that may carry fire around a selected break.

For $g\in\{T,R\}$, delete each selected non-ignition cell from $G_s^g$
for purposes of propagation, but retain $r_s$ even if selected. Write
$R_s^g(y)$ for the cells then reachable from $r_s$ and define the realized
loss by
\begin{equation}\label{eq:loss}
 Q_s^g(y)=\sum_{k\in R_s^g(y)}w_k.
\end{equation}
The optimization in this paper minimizes risk measures formed from
$Q_s^T(y)$. Evaluation on $Q_s^R(y)$ transfers the \emph{same} placement
back to the paired original graph. A separate sample of unused scenario
identifiers, if available, gives a distinct out-of-sample assessment.

\begin{proposition}[Monotonicity of graph reduction]\label{prop:reduction}
If $G_s^T$ and $G_s^R$ have the same cells, ignition, weights and treatment
convention, and $E_s^T\subseteq E_s^R$, then for every $y\in\mathcal Y$,
\begin{equation}\label{eq:graph-monotonicity}
 R_s^T(y)\subseteq R_s^R(y)
 \quad\text{and}\quad Q_s^T(y)\le Q_s^R(y).
\end{equation}
Consequently, the expected loss and every monotone risk measure, including
CVaR, of the tree-loss vector do not exceed those of the Reburn-loss vector
for this same placement.
\end{proposition}
\begin{proof}
Deleting the same selected non-ignition nodes from both graphs preserves
arc inclusion. Every remaining tree path is therefore a Reburn path.
Inclusion of the reachable sets follows, and summing nonnegative weights
gives the loss inequality. Expectation and CVaR are monotone under
componentwise loss domination.
\end{proof}

Proposition~\ref{prop:reduction} identifies the direction of the
approximation: tree propagation gives an optimistic assessment of a given
placement. It does \emph{not} imply that a placement optimal for the tree
problem is optimal for the Reburn problem, nor does it bound the difference
between those two optima without additional information. The paired
evaluation is designed to measure the discrepancy at the placements
actually selected by the tree model.
